\documentclass[10pt,a4paper]{amsart}
\usepackage[margin=19mm]{geometry}
\usepackage{amsmath,amssymb,mathtools}
\usepackage[scr=rsfs,cal=euler]{mathalfa}
\usepackage{xfrac}
\usepackage[unicode,hidelinks,bookmarks=false]{hyperref}
\newcommand{\ie}{i.e.\ }
\newcommand{\cf}{cf.\ }
\newcommand{\resp}{respectively\ }
\newcommand{\Subsection}{\subsection}
\providecommand{\nobreakdash}{\nobreak\hbox{-}}
\setlength{\emergencystretch}{2em}
\multlinegap0pt
\usepackage{amssymb}
\usepackage{mathtools}
\usepackage{xcolor}
\newtheorem{lemma}{Lemma}
\def\d{\mathrm{d}}
\title{Open problems in billiards and quantitative symplectic geometry}
\author{Bernhard Albach, Jean-Fran\c{c}ois Barraud, Misha Bialy, Johanna Bimmermann, Ana Ch\'avez C\'aliz, Mihai Damian, Lina Deschamps, Umberto Hryniewicz, Vincent Humili\`ere, Boris Khesin, Levin Maier, Agustin Moreno, Alexandru Oancea, Olga Paris-Romaskevich, Alfonso Sorrentino, Serge Tabachnikov}
\date{Source: arXiv:2602.12896v1 (CC BY 4.0)}
\begin{document}
\maketitle

\noindent\emph{Source and licence.} Open problems in billiards and quantitative symplectic geometry. Version arXiv:2602.12896v1 (CC BY 4.0), distributed by arXiv under the Creative Commons Attribution 4.0 International licence, http://creativecommons.org/licenses/by/4.0/, as recorded in the arXiv licence metadata for this exact version and shown on the abstract page https://arxiv.org/abs/2602.12896v1. Full licence notice: \texttt{licenses/arxiv-2602.12896-CC-BY-4.0.txt}.\par\smallskip

\noindent\textit{Editorial scope.} Counted unit: the article's lemma lemma1 (source0.tex lines 756--758). The article's complete original proof of it is delivered with it (source0.tex lines 760--775, one \texttt{proof} environment of 16 lines). No mathematics is written, repaired or re-derived: the statement and the proof are the article's own lines, in the article's own notation, and no retained line is substituted or re-flowed.  No further source-local context is needed: every cross-reference of every retained line resolves inside the two delivered spans.  Nothing was omitted from any retained span, so the package carries no omission record.  No retained line contains a citation, so the wrapper carries no bibliography and no bibliography entry.  No retained line contains a figure environment, an image-inclusion command or drawing code, so no image is delivered and the package declares no asset dependency.  Editorial formatting only, in addition: an AMS \texttt{amsart} preamble that restates the article's own declarations and macros used by the retained lines, namely the environments \texttt{lemma} on separate counters, together with the standard packages those lines need.  The retained environments print in this document as Lemma 1 in order of appearance, whereas the article prints the same items with the numbering of its own full document.  The complete native source of the article as distributed in the arXiv e-print for this exact version is bundled unchanged as \texttt{sources/194595/source0.tex}.
\begin{lemma}\label{lemma1}
$$ \d_\eta \circ \d_\eta =0 \qquad \Longleftrightarrow \qquad \d\eta=0.$$
\end{lemma} 

\begin{proof}
Let $\alpha \in \Omega^k(N)$. Then:
\begin{eqnarray*}
\d^2_\eta \alpha &=& \d_\eta \left(
\d\alpha - \eta \wedge \alpha
\right) = \d \left(
\d\alpha - \eta \wedge \alpha
\right) - \eta \wedge \left(
d\alpha - \eta \wedge \alpha
\right) 
\\
&=& \d^2 \alpha - \d (\eta \wedge \alpha) - \eta \wedge \d\alpha + \eta \wedge \left( \eta \wedge \alpha \right)\\
&=& \eta \wedge \d \alpha - \eta \wedge \d\alpha = 0,
\end{eqnarray*}
where in the second-last equality we used that $\d^2=0$ and $\d\eta=0$.
\end{proof}

\end{document}
